% Drafted against audit/outline.md §3, with one deviation: the outline listed
% the 2022 restatement here as well as in §2. It is covered in §2.2 and is
% cross-referenced rather than repeated. §3.2 instead explains where the
% sample boundaries come from, which the outline left implicit.
%
% Figures from data/processed/gs_filing_manifest.csv, factor_move_summary.csv
% and factor_moves.csv, hash-recorded in audit/evidence_freeze.csv at commit
% 716ecf4c.

\section{Sample and data}\label{sec:sample}

\subsection{Why one dealer over ten years}

A cross-section of banks and a single bank over a long period answer different
questions. A cross-section can establish how disclosure practice varies, which
is what \citet{perignon2010level} and \citet{campbell2022australia} do. A long
single-dealer series can ask whether one firm's reported co-movement holds up
across a change of regime, which needs the regime change inside the sample
rather than at its edge.

That is the trade this paper makes. One filer, ten consecutive fiscal years,
and enough of them that the sample contains the March 2020 dislocation, the
zero-rate period that followed, and the tightening cycle of 2022 and 2023. The
cost is that nothing here generalises to other banks: Goldman's disclosed
diversification effect is a fact about Goldman, and the stability this paper
fails to reproduce may be particular to its book. What is tested is a
reproduction, not a cross-sectional regularity.

\subsection{The sample, and where its boundaries come from}

The disclosures are the ten Form 10-K filings of The Goldman Sachs Group, Inc.,
CIK 0000886982, covering fiscal years 2016 through 2025 and filed between
February 2017 and February 2026. Appendix~\ref{sec:appendix_a} lists the
accession number, filing date and retrieval URL of each.

Both sample boundaries follow from one rule, stated in
Section~\ref{sec:disclosed}: every figure is taken from the filing that first
reported it, rather than from a later filing's restatement of it. Each 10-K
restates the prior year, so the FY2016 filing does carry 2014 and 2015 figures,
but only as comparatives. Reaching those years as first reported would
require the FY2014 and FY2015 filings, which are not in the sample because the
ten-year window was fixed before the as-first-reported rule made their absence
binding. Extending the series backwards is therefore a matter of fetching two
more filings rather than a limitation of method, and it was not done. At the recent end, FY2025 is
simply the most recent filing available at the retrieval date of 9 October
2026.

The factor panel's start date is set by the disclosed window rather than
chosen. A five-year lookback needs five years of scenarios before the first
valuation date, and the aligned trading calendar described below discards about
9\% of days, so five years of scenarios needs closer to six calendar years of
history. The panel begins on 5 January 2010 and runs to 30 December 2025,
3,670 aligned days, which leaves 1,376 days of history before the first
valuation date against the 1,260 the window consumes. One restatement in the
disclosed series is discussed in Section~\ref{sec:disclosed} and is not
repeated here; the figures used throughout are as first reported.

\subsection{The seventeen factors}

Table~\ref{tab:factors} lists the risk factors and Table~\ref{tab:factorstats}
summarises their daily moves. All seventeen come from the issuer of the
instrument, a central bank, or an exchange close: the US Treasury par yield
curve at two, five, ten and thirty years \citep{treasurycurve}; the European
Central Bank's euro area AAA central government spot curve at two, five and ten
years \citep{ecbcurve}; four equity indices; three currency pairs; and three
commodity funds. They are grouped to match Goldman's four disclosed categories
so that each sleeve can be sized against the corresponding disclosed share.

% Generated by build/make_tables.py from factor_move_summary.csv. Do not edit.

\begin{table}[tbp]
\centering
\small
\caption{The seventeen risk factors, grouped to match Goldman's four disclosed risk categories}
\label{tab:factors}
\begin{tabular}{@{}lllll@{}}
\toprule
Factor & Category & Instrument & Unit of daily move & Source \\
\midrule
us\_2y & Interest rates & Par coupon bond, semiannual & Yield change (pp) & US Treasury \\
us\_5y & Interest rates & Par coupon bond, semiannual & Yield change (pp) & US Treasury \\
us\_10y & Interest rates & Par coupon bond, semiannual & Yield change (pp) & US Treasury \\
us\_30y & Interest rates & Par coupon bond, semiannual & Yield change (pp) & US Treasury \\
eur\_2y & Interest rates & Zero-coupon bond, annual & Yield change (pp) & ECB \\
eur\_5y & Interest rates & Zero-coupon bond, annual & Yield change (pp) & ECB \\
eur\_10y & Interest rates & Zero-coupon bond, annual & Yield change (pp) & ECB \\
\addlinespace
equity\_us & Equity prices & S\&P 500 index & Fractional return & Exchange \\
equity\_europe & Equity prices & EURO STOXX 50 index & Fractional return & Exchange \\
equity\_japan & Equity prices & Nikkei 225 index & Fractional return & Exchange \\
equity\_em & Equity prices & Emerging-markets equity fund & Fractional return & Exchange \\
\addlinespace
fx\_eurusd & Currency rates & EUR/USD spot & Fractional return & Exchange \\
fx\_usdjpy & Currency rates & USD/JPY spot & Fractional return & Exchange \\
fx\_gbpusd & Currency rates & GBP/USD spot & Fractional return & Exchange \\
\addlinespace
commodity\_gold & Commodity prices & Gold fund & Fractional return & Exchange \\
commodity\_oil & Commodity prices & Crude oil fund & Fractional return & Exchange \\
commodity\_base\_metals & Commodity prices & Base metals fund & Fractional return & Exchange \\
\bottomrule
\end{tabular}
\notetext{Rate factors are repriced by full revaluation, which needs the level as well as the change, because the same 25 basis point move is worth a different amount of money at a 0.5\% yield than at a 5\% one. A yield change of 0.25 means 25 basis points throughout. Equity indices enter in local currency, so that currency risk stays in the currency sleeve rather than being counted twice. Commodities enter through exchange-traded funds rather than front-month futures; Section~\ref{sec:sample} gives the reason.}
\end{table}


% Generated by build/make_tables.py from factor_move_summary.csv. Do not edit.

\begin{table}[tbp]
\centering
\small
\caption{Daily moves by factor, 2010--2025}
\label{tab:factorstats}
\begin{tabular}{@{}lrrrr@{}}
\toprule
Factor & Mean & Std.\ dev. & Minimum & Maximum \\
\midrule
us\_2y & +0.06 & 4.73 & -57.00 & +34.00 \\
us\_5y & +0.03 & 5.62 & -35.00 & +31.00 \\
us\_10y & +0.01 & 5.63 & -36.00 & +29.00 \\
us\_30y & +0.00 & 5.51 & -45.00 & +29.00 \\
eur\_2y & +0.02 & 3.80 & -43.47 & +36.57 \\
eur\_5y & -0.01 & 4.35 & -31.25 & +29.35 \\
eur\_10y & -0.02 & 4.55 & -26.46 & +24.01 \\
\addlinespace
equity\_us & +0.056 & 1.134 & -11.984 & +9.515 \\
equity\_europe & +0.026 & 1.311 & -12.401 & +10.348 \\
equity\_japan & +0.052 & 1.381 & -12.396 & +10.226 \\
equity\_em & +0.026 & 1.409 & -12.479 & +8.053 \\
\addlinespace
fx\_eurusd & -0.004 & 0.576 & -5.001 & +3.176 \\
fx\_usdjpy & +0.016 & 0.621 & -4.552 & +4.061 \\
fx\_gbpusd & -0.003 & 0.586 & -7.604 & +3.791 \\
\addlinespace
commodity\_gold & +0.041 & 1.043 & -8.781 & +5.984 \\
commodity\_oil & -0.014 & 2.325 & -25.315 & +16.667 \\
commodity\_base\_metals & +0.014 & 1.283 & -8.371 & +6.191 \\
\bottomrule
\end{tabular}
\notetext{Equity, currency and commodity factors in percent per day; interest-rate factors in basis points per day. The two blocks are not comparable with one another, which is why the engine sizes sleeves against disclosed risk shares rather than against notional. 3,670 aligned trading days from 5 January 2010 to 30 December 2025.}
\end{table}


Seventeen factors against the roughly 70,000 the FY2016 and FY2017 filings
describe is the single largest difference between this model and Goldman's, and
the direction of its effect is known: fewer independent sources of risk means
less diversification available, which biases the measured effect downward
relative to a book with a richer factor set. It is not quantifiable here, and
Section~\ref{sec:limitations} says so rather than estimating it.

The two blocks of Table~\ref{tab:factorstats} are in different units and are
not comparable with each other. That is not a presentational awkwardness but
the reason the engine sizes sleeves against disclosed risk shares rather than
against notional: a dollar of two-year Treasury and a dollar of crude oil fund
carry risks that differ by more than an order of magnitude, so matching
Goldman's category proportions requires matching risk, not money.
Section~\ref{sec:method} gives the calibration.

\subsection{Three choices that shaped the panel}

Each of these changes the numbers, so each is stated rather than left in the
code.

\emph{Equity indices enter in local currency.} A dollar investor holding the
Nikkei carries yen risk as well as Japanese equity risk. Goldman reports
currency risk as a category of its own, so converting the indices to dollars
would push currency risk into the equity sleeve and count it twice, making the
category attribution --- the thing being matched --- meaningless.

\emph{Commodities enter through exchange-traded funds rather than front-month
futures.} On 20 April 2020 the NYMEX May 2020 WTI contract settled at
$-\$37.63$ per barrel, against a close of \$18.27 on 17 April
\citep{cftc2020wti}. A percentage return is not defined across a sign change,
and the implied one-day move of $-306\%$ is arithmetic rather than a price
movement a risk model should be asked to price. Funds cannot settle below zero.
The cost of the substitution is roll yield, which drags a fund away from spot
over long horizons; over one day it is second order, but it is a real
difference from a spot book and it is carried through to
Section~\ref{sec:limitations}.

\emph{Prices are adjusted for splits.} The crude oil fund used here ran a
one-for-eight reverse share split, effective after the close on 28 April 2020,
with post-split shares trading from the following day \citep{uso2020split}. On
unadjusted closes that is a 700\% overnight gain, and it would enter the
scenario set as a tail event in the middle of the most consequential stress
window in the sample, within days of the negative settlement above.
Split-adjusted closes are used throughout, for every traded instrument rather
than only this one.

\subsection{The trading calendar}

Tokyo, London, New York and the ECB keep different holidays, and the panel is
the intersection of their calendars rather than the union. Taking only days on
which every market traded costs about 9\% of trading days and leaves 225 to 233
observations a year against the 250 to 253 Goldman reports. The alternative,
forward-filling a closed market, would plant runs of exact zeros in the
scenario set and understate that factor's risk, which is a worse error than a
smaller sample: it would bias the thing being measured rather than merely
reducing the precision with which it is measured.

The reduced sample does cost power in every test reported in
Section~\ref{sec:validation}, and that is the price of the choice.

Two further properties of the panel are limitations rather than choices, and
neither is corrected. The surviving closes are still not simultaneous across
time zones, which depresses measured cross-market correlation and therefore
overstates diversification, in the opposite direction to the factor-count
bias above, which is why neither is used to argue the result one way or the
other. And the Treasury curve is quoted to the basis point, so at near-zero
policy rates the two-year yield shows no change at all on 21.1\% of days across
the panel, rising to 43.1\% of days in the second half of 2020. The build's
staleness check accommodates this by allowing longer runs of identical values
in yield factors than in traded prices; the longest unbroken run in the
two-year series is six days, against a limit of fifteen. The discretisation
slightly understates true rate variability through the near-zero years.
